% Copyright (c) 2026 by C. Klukas.
% Licensed under the MIT License. See LICENSE for details.

This paper has made one unit of execution carry a mobile agent: a
WebAssembly module plus one instance, whose whole state is its linear
memory and its exported mutable globals. Because the host calls into the
instance one handler at a time, it takes that memory as an image between
two calls, and the image is enough to store, copy and move the \paglet{}
(Sections~\ref{sec:model} to \ref{sec:movement}). The \paglet{} writes no
serialization code; what it gives up is strong mobility, so a long
computation is sliced by its author (Section~\ref{sec:discussion-weak}).
The rest follows from this unit: lifecycle events and failure reports
are deliveries ordered with the messages (Section~\ref{sec:lifecycle}); a
move ends in location records, so endpoints keep working while their
targets move (Section~\ref{sec:communication}); hosts share a signed
ledger and nothing else, and a \paglet{} reaches resources only through
system \paglet{}s under policy (Section~\ref{sec:mesh}); a zombie
\paglet{} keeps a replica on every host and delivers exactly one result
(Section~\ref{sec:zombies}); and the interface description in the module
lets the CLI and the web client talk to any \paglet{}
(Section~\ref{sec:tools}).

The evaluation supports these claims at the scale it was tried
(Section~\ref{sec:evaluation}). The Counter's image is 131\,212 bytes, and
one module was captured and resumed across macOS arm64, Linux arm64 and
macOS x86-64 under Rosetta~2. On the live mesh of three machines, round
trips between \paglet{}s take 1.2--5.2\,ms, content moves at
20--40\,MB/s, 100\,000 hexadecimal digits of $\pi$ take 313\,s on all
three hosts, and the Mesh Chronicle delivers its log after two takeovers
and a partition with a main on both sides. Behind the numbers stand 252
test cases run in-process and with sandboxed workers, and the check of
the zombie protocol: ten two-replica worlds explored completely, two
three-replica worlds to depth 11, 10\,000 timed runs without a failure.
What the evaluation does not support is equally plain: no move between
real machines is timed, the live numbers are single runs on one LAN, no
mesh is larger than three real or five tested hosts, and nothing was
measured with AOT or JIT, through a relay, across a NAT, or on Windows.

The next steps are the ones Sections~\ref{sec:zombies-claims} and
\ref{sec:discussion-open} name. For the zombie protocol: a checker mode
without hash compaction, a fair-cycle check for liveness, a mutation
suite, a soundness argument for the term renaming, and three replicas in
full or a cutoff argument. For the mesh: tens of hosts, relays and NAT.
For the runtime: a security review, state migration between module
versions, AOT with a per-architecture cache, suspending a handler
mid-computation, and a second guest language on the same ABI.

\paglets{}/cpp is free software under the MIT License. The sources are the
branch \code{cpp} of the repository \cite{pagletsCpp}, released as
\code{cpp-v0.2.0}; on macOS the Homebrew tap \code{cklukas/homebrew-paglets}
builds the latest release (\code{brew install cklukas/paglets/paglets});
the documentation, with the examples and the recorded results, is
published at \cite{pagletsCppDocs}, and the paper's workflow puts the
PDF of every published version beside it. The excerpts and result
tables in this paper are generated from the sources and the recorded
results by the tool that keeps the documentation current
(Appendix~\ref{app:reproduce}); a test fails whenever either falls behind
the code.
